\section*{Erd\H{o}s Problem 471}

\paragraph{Statement.}
Starting with a finite set \(Q_0\) of primes, repeatedly adjoin every
prime expressible as the sum of three distinct primes already present.
Can the sets \(Q_i\) have unbounded size?

\paragraph{External input and distinctness.}
The quantitative form of Vinogradov's three-primes theorem gives,
for every sufficiently large odd integer \(N\), at least
\(cN^2/(\log N)^3\) ordered representations \(N=p_1+p_2+p_3\),
where \(c>0\) is absolute. This standard analytic theorem is the
external dependency.
For a primary modern statement of the classical weighted formula, see
Frei, Koymans and Sofos, Section 1.2,
\url{https://arxiv.org/abs/1804.00573}.
Its singular series for odd \(N\) is bounded below by
\(2\prod_{p\geq3}(1-(p-1)^{-2})>0\).
Dividing the weighted lower bound by \((\log N)^3\), an upper bound
for each representation's weight, gives the unweighted bound used here.

Representations having two equal entries number at most \(3N\):
choose one of the three pairs of equal positions and their common
prime; the third entry is then determined. This overcounts triples
with all entries equal, which is harmless.
Since \(N^2/(\log N)^3\) grows faster than \(N\), every sufficiently
large odd integer is therefore a sum of three \emph{distinct} primes.
Every summand is smaller than \(N\).

\paragraph{Choose the seed and check the iteration.}
Choose an integer \(N_0\) past the threshold just obtained, and let
\(Q_0\) contain every prime at most \(N_0\).
We prove by induction through the ordered primes that each prime
eventually belongs to some \(Q_i\).
The assertion holds for the seed. If \(p>N_0\), write \(p=a+b+c\)
with three distinct smaller primes. By induction, they belong to
stages \(Q_{i_a},Q_{i_b},Q_{i_c}\).
Since the stages are nested, all three belong to
\(Q_{\max(i_a,i_b,i_c)}\), and the next stage contains \(p\).
This completes the induction.

Every stage is finite, because only finitely many triples can be
formed from a finite set. Their union contains all primes and is
infinite, so their cardinalities are unbounded. This is exactly the
requested existence conclusion.

\paragraph{Attribution and scope.}
The deduction from Vinogradov's theorem is credited to
Mrazovi\'c and Kova\v{c}, and independently Alon, at
\url{https://www.erdosproblems.com/471}, checked 24 September 2026.
The analytic representation bound is explicitly assumed above.
This proves existence of a finite seed; it does not settle the
additional suggested seed \(\{3,5,7,11\}\), whose success is not
needed for the existence question.

\paragraph{Completion estimate.}
The full finite-seed existence claim follows from Vinogradov's theorem.
\textbf{COMPLETION: 100\%}
